\section{Conclusion and future work}
\label{sec:conclusion}

A few unrecognised leaders can bring a flock of drones round obstacles to a
goal only they know, but how reliably depends on the flocking model more than
on anything else varied here: on whether its followers cruise, wait to be
pulled, or hold a direction of their own. A planning layer frees the leaders
from traps, and helps the flock only as far as the model lets it follow. On
the way, the project showed that a flocking model is not a controller until
the robot can fly it.

\paragraph{Future work.} Leaders that still don't communicate but share a
convention, such as passing every obstacle on the same side, which would
remove the split that breaks the hybrid without a leader having to know the
others; leaders that check they are followed and
wait or replan; a stopping limit shared by two neighbours, to keep crowded
drones apart; flight over obstacles in 3D; obstacles sensed rather than
known, and real Crazyflies; more seeds and maps for the smaller effects.
